\clearpage\begingroup\SourceLeaf{115b}\HistoricalHead{110}{L}{PHYSIK. BRIEFWECHSEL.}
\begin{center}[24.]\end{center}
\begin{center}\textsc{Gauss} an \textsc{Olbers}.\quad Göttingen, 23. Juli 1836.\end{center}
\begin{center}\begin{minipage}{.88\linewidth}\centering\hrule\vspace{3pt}
\small [\textsc{Wilhelm Olbers}, Sein Leben und seine Werke II, 2, S. 642.]
\par\vspace{3pt}\hrule\end{minipage}\end{center}
\noindent .\ .\ .\ .\ .\ .\ .\ .\ .\ .\ .\ .\ .\ .\ .\ .\ .\ .\ .\ .\ .\par
\indent Unter meinen jetzigen magnetischen Experimenten ist das merkwürdigste
eines, wo die Inductionswirkung des Erdmagnetismus auf einen 700fachen
Drahtring, übersponnen und zusammen über 13000 Fuss lang, bestimmt wird.
Dieser Ring oder dieses Rad wird um eine Axe gedrehet, die einen horizontalen
Diameter des Rades bildet. Diese Drehungsaxe macht genau einen rechten
Winkel mit dem magnetischen Meridian, und die beiden Enden des Drahts
sind bis zum Multiplicator des grossen Magnetometers der Sternwarte fortgeführt,
welcher Multiplicator \textls[120]{jetzt} aus 610 Umwindungen übersponnenen
Drahts besteht. Die Drehung geschieht tactmässig nach der Uhr, alle 2 Secunden
Eine Umdrehung, und nach jeder halben Umdrehung wechselt vermöge eines
eigenthümlichen Mechanismus die Verbindung der Drahtendungen des Rades mit
ihren Fortsetzungen. Die Einrichtung ist nach Vorschrift der Theorie so, dass
dieser Wechsel genau Statt findet, wenn das Rad dem magnetischen Aequator
parallel ist, oder einfacher gesagt, normal gegen die Richtung der Erdmagnetischen
Kraft. So wirkt der durch die Induction erzeugte Strom immer in
Einerlei Sinn auf das Magnetometer und bewirkt, obgleich die ganze Drahtlänge
so gegen 20000 Fuss lang ist, doch noch Ausweichung von mehreren
hundert Scalentheilen an der 25pfündigen Nadel. Bei einigen Versuchen
wird auch noch der Draht des Inductors (\textsc{Schumachers} J. B. 1836, S. 41\textsuperscript{[*)]}),
\begingroup\renewcommand{\thefootnote}{}\footnotetext{[*) \textit{Erdmagnetismus und Magnetometer}, Jahrbuch für 1836, herausgegeben von \textsc{Schumacher}, Werke V, S. 315, siehe insbesondere S. 340.]}\endgroup
\clearpage\SourceLeaf{116}\HistoricalHead{111}{R}{GAUSS AN GERLING UND OLBERS.}
\noindent der aber jetzt nicht mehr 3537, sondern 7000 Umwindungen hat, mit in die
Kette aufgenommen, wo also der Strom gegen 27000 pariser Fuss sehr dünnen
Drahts zu durchlaufen hat und dann noch immer 325 Scalentheile Ausschlag
gibt. Diese Versuche machen einen Theil von denjenigen aus, die zum
Zweck haben, alles was sich auf die Wechselwirkung zwischen Magnetismus
und galvanischen Strömen bezieht, auf absolute Maasse zu bringen. Genäherte
Zahlenbestimmungen geben schon meine früheren Versuche. Die gegenwärtigen
werden aber die Genauigkeit noch sehr vergrössern.\quad .\ .\ .\ .\ .\ .\ .\ .\ .\ .\ .\ .\ .\ .

\bigskip\Needspace{13\baselineskip}
\begin{center}[25.]\end{center}
\begin{center}\textsc{Gauss} an \textsc{Gerling}.\quad Göttingen, 19. December 1836.\end{center}
\noindent .\ .\ .\ .\ .\ .\ .\ .\ .\ .\ .\ .\ .\ .\ .\ .\ .\ .\ .\ .\par
\indent Mein Sohn ist jetzt wieder in Stade; er hat sich 3 Wochen in Paris
aufgehalten und hat von den dortigen Gelehrten viel Gefälligkeit erfahren,
besonders von \textsc{Poisson} und \textsc{Libri}. Auch mit \textsc{Arago} ist er bekannt geworden,
welcher ihm erklärt hat, er ziehe doch die \textsc{Gambey}schen Apparate den
meinigen vor, weil letztere \textls[120]{zu viel} Anomalien zeigten. \textsc{Airy-Christis} Bedenklichkeiten
bestanden hingegen, wie Sie sich erinnern, darin, dass sie
meinten, meine Apparate zeigten \textls[120]{zu wenig} Anomalien.\quad .\ .\ .\ .\ .\ .\ .\ .\ .\ .\ .\ .
\clearpage\endgroup
